% !PRINTED 265
\weberpage{265}

\begin{center}
\textbf{\S\,74. Die \gesperrt{Form} der \gesperrt{Schlaeflischen} Modulargleichungen}\\
\textbf{f\"ur einen Primzahlgrad.}
\end{center}

Die Form, die wir im vorigen Paragraphen f\"ur die zwischen
\weq{1}{u=f(\omega),\qquad v=f\Bigl(\frac{c+\partial\omega}{a}\Bigr)}
% Source confirms the printed partial-shaped d glyph, transcribed as \partial.
bestehenden Relationen gefunden haben, l\"a\ss t sich, wenigstens
wenn der Transformationsgrad eine Primzahl $p$ ist, leicht unter
ein allgemeines Gesetz bringen. Da f\"ur die Folge viel auf diese
Form ankommt, gehen wir hier noch etwas genauer darauf ein.

Die Bestimmung der Zahlen $r, s$ nach (11) und (15) des
vorigen Paragraphen h\"angt von dem Verhalten von $p$ gegen den
Modul 24 ab, und da wir den Fall $p=3$ ausschlie\ss en k\"onnen,
so haben wir folgende F\"alle:
\weq{2}{\begin{array}{lll}
p\equiv 1 \ (\mathrm{mod}\ 24) & r=1 & s=12\\
p\equiv 5 & r=3 & s=2\\
p\equiv 7 & r=4 & s=3\\
p\equiv 11 & r=6 & s=1\\
p\equiv 13 & r=1 & s=6\\
p\equiv 17 & r=3 & s=4\\
p\equiv 19 & r=2 & s=3\\
p\equiv 23 & r=12 & s=1,
\end{array}}
so da\ss\ $r+s$ stets ungerade ist und $p+1$ durch $2r$, $p-1$
durch $2s$ teilbar ist. Hiernach wird
\weq{3}{\begin{aligned}
A &= \Bigl(\frac{u}{v}\Bigr)^r+\Bigl(\frac{v}{u}\Bigr)^r,\\
B &= (uv)^s+\Bigl(\frac{2}{p}\Bigr)\frac{2^s}{(uv)^s}.
\end{aligned}}
Nun wissen wir, da\ss\ die $p+1$ Gr\"o\ss en $v$ die Wurzeln einer
irreducibeln Transformationsgleichung
\weq{4}{\Phi_p(v,u)=v^{p+1}+U_1v^p+\cdots U_{p+1}=0}
sind, in der die Koeffizienten $U_1\ldots, U_{p+1}$ rationale
Funktionen von $u$ sind.

Wir schlie\ss en sofort, da\ss\ es \gesperrt{ganze} rationale Funktionen
von $u$ sind. Denn \gesperrt{erstens} werden f\"ur $u=\infty$ die s\"amtlichen
Wurzeln von (4) unendlich, wie man erkennt, wenn man $\omega=i\infty$,
also $q=0$ werden l\"a\ss t. Zweitens geht nach (9) des vorigen
%CONT>>
